\documentclass[11pt]{article}
\usepackage[margin=1in]{geometry}
\usepackage[T1]{fontenc}
\usepackage{lmodern}
\usepackage{amsmath,amssymb,amsthm}
\usepackage[hidelinks]{hyperref}
\usepackage{microtype}
\usepackage{enumitem}
\raggedright

\newtheorem*{theorem}{Theorem}
\newtheorem*{lemma}{Lemma}

\newcommand{\B}{\mathcal B}
\newcommand{\R}{\mathbb R}

\begin{document}

\begin{center}
{\Large Full solution packet for arXiv:1810.05356}\\[4pt]
{\large The Fatou property for Bochner-integrable Banach-lattice-valued functions}
\end{center}

\paragraph{Status.}
\textbf{full\_solution\_likely\_valid}.  This packet answers the question
following Theorem~50 in Omid Zabeti, \emph{On the lattice structure of the
space of all Bochner integrable Banach lattice-valued functions},
arXiv:1810.05356.

\paragraph{Source question.}
The source proves that, for a finite measure space $(X,\Sigma,\mu)$ and a
Banach lattice $E$, the Bochner lattice $\B(X,E,\mu)$ has the sequential Fatou
property if and only if $E$ has the sequential Fatou property.  Immediately
after Theorem~50, it asks whether the same equivalence remains true for the
full Fatou property, i.e. for positive nets instead of positive sequences.

\paragraph{Convention.}
We assume $\mu(X)>0$.  This is the nondegenerate case in which the source's
constant-function embedding $E\to \B(X,E,\mu)$ is a topological lattice
embedding.  If $\mu(X)=0$, then $\B(X,E,\mu)=\{0\}$ is Fatou for every $E$, so
the converse direction is false without this harmless convention.

\paragraph{Definitions.}
The norm of a Banach lattice $E$ is Fatou if whenever $0\leq x_\alpha\uparrow x$
in $E$, one has
\[
  \|x\|=\sup_\alpha \|x_\alpha\|.
\]
We write $\B(X,E,\mu)=L^1(\mu;E)$ for the Banach lattice of Bochner-integrable
$E$-valued functions modulo a.e. equality, with norm
\[
  \|f\|_{\B}=\int_X \|f(w)\|\,d\mu(w).
\]

\paragraph{Proof Intuition.}
The sequential proof in the source uses pointwise monotone convergence.  For
nets, the same idea still works once one uses the standard order structure of
Bochner $L^1$: an increasing net with a Bochner-integrable supremum is
represented by an a.e. pointwise increasing family.  At almost every point,
Fatou in $E$ converts the vector supremum into the scalar supremum of the
norms.  Then the scalar Beppo Levi theorem for upward directed families
passes the supremum through the integral.

\begin{lemma}[Scalar directed monotone convergence]
Let $(h_\alpha)$ be an upward directed family of nonnegative measurable scalar
functions on a finite measure space, and let $h=\sup_\alpha h_\alpha$ be
measurable.  Then
\[
  \int_X h\,d\mu=\sup_\alpha \int_X h_\alpha\,d\mu .
\]
\end{lemma}

\begin{proof}
This is the usual monotone convergence theorem in directed-family form.  The
point is that the supremum over any finite subfamily is again dominated by
one member of the directed family, so the standard proof for increasing
sequences applies to the net of finite suprema.
\end{proof}

\begin{lemma}[Order representation used]
In the Bochner lattice $L^1(\mu;E)$, if $0\leq f_\alpha\uparrow f$, then the
classes admit representatives, after changing on a null set, such that
$0\leq f_\alpha(w)\uparrow f(w)$ in $E$ for almost every $w$.
\end{lemma}

\begin{proof}
This is the standard pointwise representation of lattice suprema in
Bochner-integrable Banach-lattice-valued function spaces; it is the same
lattice structure used in the source via van Zuijlen's theorem that
$\B(X,E,\mu)$ is a Banach lattice under pointwise a.e. operations.
\end{proof}

\begin{theorem}
Let $(X,\Sigma,\mu)$ be a finite measure space with $\mu(X)>0$, and let $E$ be
a Banach lattice.  Then $E$ has the Fatou property if and only if
$\B(X,E,\mu)=L^1(\mu;E)$ has the Fatou property.
\end{theorem}

\begin{proof}
Assume first that $E$ has the Fatou property.  Let
$0\leq f_\alpha\uparrow f$ in $\B(X,E,\mu)$.  By the order representation
lemma, choose representatives such that
$0\leq f_\alpha(w)\uparrow f(w)$ in $E$ for almost every $w$.  For such $w$,
the Fatou property of $E$ gives
\[
  \|f(w)\|=\sup_\alpha \|f_\alpha(w)\|.
\]
The scalar functions $h_\alpha(w)=\|f_\alpha(w)\|$ form an upward directed
family of nonnegative measurable functions with measurable supremum
$h(w)=\|f(w)\|$.  Hence scalar directed monotone convergence yields
\[
\begin{aligned}
  \|f\|_{\B}
  &= \int_X \|f(w)\|\,d\mu(w) \\
  &= \int_X \sup_\alpha \|f_\alpha(w)\|\,d\mu(w) \\
  &= \sup_\alpha \int_X \|f_\alpha(w)\|\,d\mu(w)
   = \sup_\alpha \|f_\alpha\|_{\B}.
\end{aligned}
\]
Thus $\B(X,E,\mu)$ is Fatou.

Conversely, assume that $\B(X,E,\mu)$ is Fatou.  Let
$0\leq x_\alpha\uparrow x$ in $E$.  For each $\alpha$, let
$c_\alpha(w)=x_\alpha$ and $c(w)=x$ be the corresponding constant Bochner
functions.  Then $0\leq c_\alpha\uparrow c$ in $\B(X,E,\mu)$, so
\[
  \mu(X)\|x\|=\|c\|_{\B}
  =\sup_\alpha \|c_\alpha\|_{\B}
  =\mu(X)\sup_\alpha\|x_\alpha\|.
\]
Since $\mu(X)>0$, division by $\mu(X)$ gives
$\|x\|=\sup_\alpha\|x_\alpha\|$.  Hence $E$ is Fatou.
\end{proof}

\paragraph{Verification notes.}
The proof is exactly the source's sequential argument with two net-level
replacements: pointwise a.e. representation of order suprema in Bochner
$L^1$, and scalar monotone convergence for upward directed families.  The
zero-measure edge case is explicitly separated because the source's Lemma~60
also needs positive total measure to embed $E$ nontrivially by constant
functions.

\paragraph{Novelty check.}
The run indexes had no direct entry for arXiv:1810.05356.  Web and local
searches used phrases including ``Bochner Fatou property Banach lattice
L1(E)'', ``Fatou property Bochner integrable Banach lattice'', and
``Zabeti Bochner integrable Fatou''.  The bounded search found the source
paper but no later explicit answer to the question.

\paragraph{Human-review recommendation.}
Review as a short affirmative solution.  The only substantive background step
to verify is the standard order-representation lemma for arbitrary increasing
nets in Bochner $L^1$; if a reviewer wants the packet fully self-contained,
that lemma can be expanded from the usual construction of lattice suprema in
spaces of a.e. equivalence classes of Bochner-measurable Banach-lattice-valued
functions.

\begin{thebibliography}{9}
\bibitem{Zabeti1810}
O. Zabeti,
\emph{On the lattice structure of the space of all Bochner integrable Banach
lattice-valued functions},
arXiv:1810.05356.

\bibitem{VanZuijlen2012}
W. van Zuijlen,
\emph{Integration of functions with values in a Riesz space},
Master thesis, Radboud Universiteit Nijmegen, 2012.
\end{thebibliography}

\end{document}
